\chapter{ASSALTO}\label{assalto}

\hfill\break

Giraud mi chiamò, aveva bisogno di parlare del piano d'assalto alla base
sull'asteroide. Lo trovai, insieme al sergente Thompson, in una stiva di
carico vicino alla palestra. Skippy mi aveva detto che stava elaborando
un piano d'assalto con Giraud e io avevo lasciato che lavorasse senza
inutili interferenze da parte mia.

«Buongiorno, signore», disse il sergente Thompson, con il suo
affascinante accento inglese.

«Buongiorno», risposi, con ancora in bocca il sapore della melma che
avevo mangiato a colazione. Quella mattina, avevo provato a mescolare la
melma mela-cannella, che non aveva avuto successo, con una melma al
cioccolato che non era troppo male, il risultato fu qualcosa che ero
stato felice di ingoiare il più in fretta possibile, in modo da non
doverne sentire il sapore. C'era una tacita competizione tra i membri
dell'equipaggio per trovare combinazioni di buon sapore, che venivano
poi pubblicate nel nostro network zPhone interno. Il mio esperimento di
quella mattina stava per finire nella colonna ``Falliti''. «Tenente,
vedo che sei stato occupato.»

«Sì, signore», annuì Giraud. Quando si era arruolato nella missione,
avevo avuto l'impressione che mi considerasse un po' uno scemo
fortunato, qualcuno che aveva le chiavi del futuro, ma non sapeva come
usarle e aveva bisogno di soldati professionisti come lui per prendere
le decisioni importanti. È probabile che avesse ancora un'impressione
meno che lusinghiera delle mie abilità di pianificazione tattica, ma
nutriva un riluttante rispetto per il modo in cui avevo gestito in
generale la missione fino a quel momento. La mia idea di liberarmi del
problema kristang, facendo saltare le loro astronavi in un gigante
gassoso, aveva impressionato molte persone. Questo, e il fatto che la
missione avesse avuto un enorme successo fino a quel momento e fosse
andata proprio come avevo promesso. Eravamo fuggiti da Paradiso, avevamo
preso una fregata kristang e una nave madre thuranin e ora stavamo
saltando verso un wormhole. Il successo incoraggia la fiducia. Avevo
bisogno di continuare a mietere successi. «Il piano d'assalto alla base
sull'asteroide richiede l'elemento sorpresa; dobbiamo muoverci in fretta
una volta che le navi sono nell'hangar d'atterraggio. Ciò significa che
i componenti della squadra d'assalto non possono aspettare che le porte
dell'hangar si chiudano e che l'hangar si ripressurizzi, devono uscire
dalle navi e accedere all'interno della base prima che i kristang
possano reagire. Per farlo, hanno bisogno di tute spaziali.»

Giraud indicò con un gesto uno scaffale su cui c'erano due tute
spaziali, una piccola che sembrava leggera e una più grande, che era
ingombrante e rivestita di una pesante corazza. «Questa tuta thuranin»,
indicò quella piccola, «non andrebbe bene neppure per il membro
dell'equipaggio più basso. Skippy mi ha detto che non c'è modo di
espandere le dimensioni di queste tute, le strutture di fabbricazione di
questa nave non hanno la capacità di produrre i materiali.»

«Vero», concordò Skippy dall'altoparlante sulla paratia. «I gruppi
tattici thuranin includono navi di supporto con strutture di
fabbricazione complete, le singole navi da guerra hanno solo capacità di
fabbricazione limitate.»

«Questo era un problema serio», continuò Giraud, «finché il sergente
Thompson non ha pensato alle tute spaziali sulla nostra nave kristang.
Quelle tute sono corazzate e potenziate per il combattimento, i motori
negli arti migliorano i movimenti di chi le indossa, in modo che il peso
non sia un problema per gli esseri umani.»

Le tute kristang sembravano fantastiche, le raggiunsi e battei le nocche
sull'armatura. Era solida. Le armi potevano essere attaccate alle staffe
ai polsi, quindi chi le indossava doveva solo controllare il grilletto.
La visiera era di una sorta di vetro affumicato, forse non era neanche
vetro. Mi fecero venire in mente Iron Man, o una tuta meccanica da Halo,
o una dozzina di altri videogiochi. «Buona idea, sergente. La domanda
successiva è: queste tute non sono troppo grandi?»

«Questo è il problema», ammise Giraud.

Thompson girò attorno alla tuta kristang, si fermò sul retro, dove la
schiena si apriva, e s'infilò dentro con cautela. Giraud lo aiutò a
tirarla su e a chiuderla sul retro. Gli indicatori di stato si accesero
sui display del polso e la tuta fece due cauti passi avanti, poi il
pannello frontale si ritrasse in modo da mostrare il volto di Thompson.
«Non è l'ideale», disse il sergente, le sue parole uscirono distorte. Il
suo volto era nascosto dal naso in giù. «Il mio mento colpisce il fondo
del casco, non riesco ad aprire la bocca per parlare bene.»

«Il sergente Thompson è alto quasi due metri», disse Giraud.

«Sono sei piedi e quattro per te, Joe», aggiunse Skippy per dare una
mano.

«Ed è a malapena in grado di utilizzare la tuta in modo efficace»,
continuò Giraud. «La tuta è regolabile e Skippy dice che può fabbricare
alcuni componenti per noi.»

«Un numero limitato di componenti», ammonì Skippy, «non vi gasate
troppo.»

«Quanto regolabile?» Mi alzai sulle dita dei piedi per guardare nel
casco. Thompson non sembrava essere a suo agio.

«La regolazione massima, per l'uso pratico in combattimento, copre fino
a un metro e ottantadue di altezza, cioè sei piedi, non al di sotto»,
precisò Skippy.

«Io sono sei e tre», meditai, «chi altro di noi?» Thompson e Giraud
erano alti, così come Chang.

«Otto persone», Giraud contò sulle dita, «io, lei, il sergente Thompson,
il colonnello Chang, lo specialista Putri e i soldati semplici Marsden,
Darzi e Putri.»

«Otto persone? Bastano per la riuscita di un assalto?», chiesi.

«Sette persone», mi corresse Giraud. «Lei non parteciperà all'assalto,
signore. Deve restare a bordo dell'Olandese e dirigere l'operazione.»

«Ha ragione, Joe», aggiunse Skippy. «Io starò qui a guidare le navicelle
e monitorare i sensori, c'è bisogno di qualcuno che mi dica cosa fare.»

Cercando di non lasciar trapelare l'irritazione sulla mia faccia, dissi
brusco: «Ne discuteremo più tardi. Sette persone bastano perché la
missione riesca?».

«No», dichiarò Giraud in tono monocorde.

«Sono d'accordo, Joe», disse Skippy. «È qui che ci serve il tuo genio
militare. Potrei spaccare l'asteroide con dei cannoni a rotaia e poi
cercare tra i detriti, ma il rischio di danneggiare gli oggetti di cui
abbiamo bisogno è troppo alto.»

«Tenente», mi rivolsi a Giraud, «tu hai esperienza e addestramento nelle
forze speciali. Se pensi che non possiamo riuscire nella missione,
allora non ho niente da aggiungere.»

Giraud si accigliò: «Siamo tornati al via».

«Non ci arrenderemo. Sono sicuro che hai considerato tutto. Facciamo un
passo indietro, facciamo una pausa e affrontiamo la cosa di nuovo domani
a mente fresca.»

Giraud aiutò Thompson a uscire dalla tuta da combattimento kristang,
Thompson mosse le spalle irrigidite, la tuta gli stava a stento e i
bordi gli avevano scavato a fondo sotto le braccia.

«Sergente», chiesi, mentre esaminavo la tuta spaziale thuranin, «se
quella tuta kristang è regolata nel modo giusto, sei sicuro di poterla
usare in combattimento?»

«Sì, signore. Quando si muove, non si fa quasi fatica. Con la pratica,
possiamo essere letali con queste tute. Siamo fortunati che i kristang
non siano riusciti a usarle quando abbiamo preso la Fiore.»

«Bene, bene.» Ero distratto. La tuta thuranin era leggera, quasi
inconsistente. Anche supponendo che fosse fatta di materiale esotico ad
alta tecnologia, non riuscivo a immaginare come potesse competere in
combattimento contro una tuta kristang. «Skippy, queste tute thuranin mi
ricordano lo spandex. Questo materiale si indurisce quanto un diamante o
qualcosa del genere? O hanno, tipo, scudi energetici come una nave?»

«No, niente di così raffinato. È una tuta spaziale», rispose Skippy.

«Ehm. Come fanno i thuranin a combattere dentro queste cose?»

«Combattere?» Skippy rise. «I thuranin di solito non vanno in
combattimento. Quelle piccole zucche vuote verdi non si esporranno in
modo diretto al pericolo. Controllano da remoto i droni militari, è
probabile che tu li chiameresti droidi da combattimento o robot o
qualcosa del genere.»

«Droidi da combattimento? Ci sono droidi da combattimento su questa
nave?»

«Certo, ce ne sono tre dozzine nell'armeria della nave», c'era un ``Ma
tu pensa'' sottinteso nella voce di Skippy, «trentotto unità, per essere
precisi.»

«Armeria? Quale armeria?», chiesi.

«Abbiamo robot da combattimento?», si inserì Giraud. «Sarebbe stata una
bella cosa saperlo prima di trascorrere giorni a pianificare un
assalto!»

«Ci sono robot da combattimento, fucili, razzi, giocattoli di tutti i
tipi nell'armeria, si trova a poppa della sezione di comando. È un po'
nascosta, bisogna sapere dove guardare.»

«Perché cazzo non ce l'hai detto?» Non cercai d'impedire alla
frustrazione di trapelare dalla mia voce.

«Beh, Joe, non l'hai chiesto», mormorò Skippy, «e non mi offrirò
volontario per dire a un branco di scimmie dove possono trovare
giocattoli pericolosi con cui trastullarsi.»

«Merde.» Giraud era paonazzo. «Colonnello, se mai avesse voglia di
buttare Skippy fuori da una camera d'equilibrio, me lo faccia sapere e
sarò molto felice di occuparmene per lei.»

«Cosa?», chiese Skippy con fare innocente. «Cosa ho fatto?»

Alzai il pollice per Giraud. «Skippy, se non lo capisci da solo, non
posso stare qui a spiegartelo io. Questi robot da combattimento, gli
umani sono in grado di controllarli, oppure è una cosa che i thuranin
fanno attraverso la loro cibernetica?»

«I thuranin usano la cibernetica, certo, che è una delle modalità più
efficienti di manipolare dispositivi da remoto, soprattutto in
combattimento, dove il tempo di risposta è fondamentale. Niente ci
impedisce di collegare qualcosa che consenta agli umani di fungere da
agenti di manipolazione da remoto.»

Sarebbe stato d'aiuto per me se Skippy si fosse dato la pena di spiegare
termini tecnici inconsueti. «Per chiarire, ``manipolazione da remoto''
significa controllo a distanza? Con che cosa, joystick, pulsanti, cose
così?»

«Non ci posso credere! No, amico, sarebbe troppo lento! Cos'è, pensi di
essere nel 1985 e di giocare a Super Mario? Attaccheremo dei sensori
agli operatori, così i robot si muoveranno mentre si muovono loro.
Mocap, cattura del movimento, molto meglio rispetto alla tecnologia
usata da voi scimmie. Ora che ci penso, ci servono anche gli occhiali,
così l'operatore può vedere quello che vede il robot. Non è così
efficiente come agganciarsi in modo diretto al nervo ottico come fanno i
thuranin, ma dovrebbe andare abbastanza bene. E i sensori dovrebbero
includere gli accelerometri, in modo che l'operatore possa ottenere un
feedback quando il robot incontra resistenza.»

«Saresti così gentile da aprirci l'armeria, in modo che noi sporche
scimmie possiamo vedere questi robot incredibili?» Misi quanto più
sarcasmo possibile nella mia voce.

«Visto che l'hai chiesto in modo così gentile, certo.» Skippy non
resisteva alla tentazione di lanciarci frecciatine: «Prima lavati le
zampe, però, per favore».

\hfill\break

I droidi erano impressionanti, anche se il sergente Adams coniò in
fretta il termine ``combot''²⁸ per definirli. Secondo lei, non erano
veri androidi, perché avevano solo una limitata capacità di movimento e
reazione in autonomia e richiedevano il controllo di un operatore
senziente. Fatti uscire tre combot dall'armeria ed entrati in una stiva
di carico vuota, Giraud, Adams e Thompson furono in grado di
controllarli senz'alcuna attrezzatura speciale, bastava che Skippy
guardasse i movimenti degli umani e istruisse i combot a seguirli. La
nostra lattina di birra super-intelligente ci ammonì che in
combattimento non avremmo potuto fare affidamento su di lui, la base
sull'asteroide aveva una pesante schermatura e gli operatori dovevano
essere vicini al loro combot, per ridurre il ritardo temporale. Durante
l'assalto, dovevano quindi essere a bordo delle navi nell'hangar
d'atterraggio. Non ne ero felice, avremmo messo più persone a rischio,
anche se il rischio per loro sarebbe stato inferiore a quello che
avrebbero corso con indosso armature potenziate kristang. Avevamo ancora
bisogno di soldati in tuta spaziale nell'assalto, Skippy non si fidava a
delegare ai combot compiti che richiedessero raffinate abilità motorie,
come selezionare, raccogliere e trasportare la fragile, antica
attrezzatura degli Anziani di cui aveva bisogno. Inoltre, utilizzando
soltanto i combot avremmo puntato tutto su una carta sola: secondo i
dati di kristang e thuranin che Skippy aveva scaricato, tuttavia, i
primi avevano progettato le loro difese nello specifico per proteggere
la base dai loro patroni, e non c'erano garanzie del fatto che le
lucertole non avessero la capacità di interferire con la connessione
della manipolazione da remoto. I kristang sapevano tutto sui combot
thuranin e dovevano avere speso molto tempo e fatica per capire come
sconfiggerli.

Il piano rivisto di Giraud prevedeva sei persone con indosso armature
kristang e nove operatori di manipolazione da remoto nelle navicelle.
Skippy avrebbe pilotato a distanza queste ultime, il che significava che
non avremmo avuto bisogno d'impiegare persone nel ruolo di piloti, ma io
volevo comunque che almeno quattro dei miei fossero addestrate a
pilotare navicelle thuranin. Sei persone avevano bisogno che le tute
kristang fossero regolate per adattarsi a loro, poi di esercitarsi,
indossandole nelle simulazioni di combattimento. Quattro avevano bisogno
di formazione per pilotare navicelle thuranin, e sarebbe stata una
formazione minima, perché gli stessi, insieme a quasi tutti gli altri,
avrebbero dovuto esercitarsi nella manipolazione da remoto dei combot.
Con riluttanza, concordai con Giraud, Chang e Simms che durante
l'assalto io sarei rimasto a bordo dell'Olandese insieme a Skippy, Desai
e al soldato semplice Walorski.

Walorski era stato svegliato e rilasciato dalla capsula medica da
Skippy, dopo aver deciso che il suo stato di salute era migliorato a
sufficienza da affrontare il resto del regime di guarigione in
autonomia. In autonomia, non fosse che il suo avambraccio sinistro, che
era stato frantumato e quasi reciso durante la battaglia per prendere la
Fiore, era racchiuso in una manica rigida progettata per le gambe dei
thuranin. Era un po' troppo grande per l'avambraccio di Walorski, tanto
che continuava a sbattere contro le cose, causando dolore
all'infortunato e una frustrazione infinita al dottor Skippy. Ogni
piccolo trauma all'avambraccio in via di guarigione rallentava un po' il
recupero di Walorski. Il manicotto conteneva nano-sonde che rinsaldavano
le ossa, i nervi e la muscolatura e fiale di fluidi che dovevano essere
cambiate tre volte al giorno. Skippy predisse che la completa
guarigione, con rimozione del manicotto, sarebbe avvenuta nel giro di
tre settimane, forse prima, se Walorski avesse collaborato. Certo una
prognosi decisamente migliore di quella che temevamo in origine:
l'amputazione del braccio sinistro sotto il gomito e un'alta probabilità
di morire, date le scarse scorte mediche che avevamo portato con noi e
la mancanza di un medico umano.

\hfill\break

Si potrebbe pensare che Walorski fosse grato, e lo era, ma era anche un
soldato e vedeva tutti gli altri, tranne me e Desai, esercitarsi con
l'armatura kristang o con i combot. Protestò: gli doveva essere affidato
un combot, avrebbe dovuto fare strada con una navetta durante
l'incursione, non era qualificato come pilota e non voleva restare nelle
retrovie mentre tutti gli altri rischiavano le loro vite in una
battaglia cruciale per la sopravvivenza dell'umanità. Lo disse a me, il
suo comandante, che sarebbe rimasto sull'Olandese mentre tutti gli altri
avrebbero rischiato la vita in una battaglia fondamentale per la
sopravvivenza dell'umanità. Walorski doveva lavorare sulle sue abilità
d'interazione umana, questo era certo. Mi sarei sentito anch'io così,
riluttante all'idea di essere lasciato indietro mentre i miei compagni
soldati combattevano contro il nemico? Cazzo, sì, è così che mi sentivo.
È probabile che l'assalto alla base sull'asteroide avrebbe deciso se il
nostro pianeta natale, la nostra intera specie sarebbero rimasti o meno
sotto lo stivale dei kristang; avrebbe potuto decidere se l'umanità
sarebbe sopravvissuta in qualche forma riconoscibile. Non avrei mai
voluto rimanere indietro in relativa sicurezza sull'Olandese, in grado
di saltare via se la nave fosse stata minacciata. Dopo la battaglia,
qualunque cosa fosse successa, avrei sempre saputo che la mia parte
nell'azione era stata quella di starmene seduto su una sedia. E anche
tutti gli altri lo sapevano. I miei tentativi di convincere il soldato
Walorski che era necessario a bordo dell'Olandese come pilota di riserva
erano ostacolati dal fatto che non ne ero del tutto convinto io stesso.
Walorski era sul punto di sconfinare nell'aperta insubordinazione,
quando Skippy intervenne affermando che non avrebbe potuto essere utile
come controller di un combot senza la piena funzionalità della mano
sinistra. Walorski avrebbe potuto aiutare Desai a pilotare la nave con
una mano o avrebbe potuto starsene seduto senza fare niente, ma durante
l'assalto sarebbe stato solo d'impaccio. Era una gara a chi era più
infelice: Walorski o io.

\hfill\break

C'era un ampio gruppo di persone infelici a bordo della nave: tutti
quelli dediti all'addestramento con armature potenziate o combot nella
stiva di carico che avevamo preparato per le esercitazioni di
combattimento. Erano scontenti perché Giraud aveva commesso l'errore di
assegnare al sergente Adams il compito di convertire una stiva di carico
in un simulatore di combattimento e di progettare il programma di
addestramento. Nessuna delle due cose era un problema. Mentre dirigeva
l'addestramento, Adams sceglieva la musica da far risuonare a tutto
volume dagli altoparlanti nella stiva di carico. La musica alta
deconcentrava, contribuiva a addestrare le persone a lavorare in
presenza di distrazioni, con le orecchie assalite da rumori così forti
da non riuscire a pensare con lucidità. Come accadeva in combattimento.
Il problema non era la musica ad alto volume di per sé, ma i gusti
atroci di Adams in merito. Quando sto svolgendo un'attività fisica
intensa, che sia per addestrarmi al combattimento oppure semplicemente
un allenamento in palestra, voglio delle melodie che facciano pompare il
cuore, che facciano stringere i pugni nell'esultanza. Rock`n'roll, rap,
qualcosa con un bel ritmo.

Quella che Adams faceva riprodurre dalla sua playlist era musica gospel,
cajun, jazz, polka e, non sto scherzando, bluegrass. Erba blu del cazzo,
tipo banjo e ragazzi che cantano col naso, con quella piagnucolosa
affettazione che mi fa venire voglia di spaccare loro il banjo sulla
testa. Non sono un grande fan del bluegrass, nel caso ve lo steste
chiedendo. Country sì, bluegrass no. E nessuna delle canzoni era un buon
esempio del suo genere, questo è certo, Adams aveva scavato per trovare
quel genere di composizioni di merda che gli artisti usano per riempire
i loro album quando sono a corto di idee, o talento, o entrambi. Ah, e
faceva girare la stessa orribile canzone ancora e ancora e ancora e
ancora e ancora. Prova a concentrarti mentre quella roba ti fa esplodere
le orecchie.

Ne deriva che Adams era un genio del male. Se ti piace la musica che
stai ascoltando, puoi prendere il ritmo e non ti deconcentra. Se la
musica ti giunge alle orecchie come il rumore di unghie strisciate su
una lavagna, dedicherai parte della tua facoltà di pensiero a
fantasticare sul modo migliore per uccidere la persona che ha scelto la
playlist e sarai distratto, e lo scopo dell'addestramento è quello di
imparare a concentrarsi e ignorare le distrazioni. Il metodo di Adams
era molto efficace.

Ma la musica faceva comunque schifo. La prima volta che mi trovai a
guardare un esercizio, non appena fu finito, rispolverai la vecchia,
vecchia scuola e misi su un po' di Coolio. Questo mi fece guadagnare una
spolliciata in su di sollievo da parte dell'equipaggio: era stato il
miglior incitamento morale ricevuto da quando avevano visto le navi
kristang indesiderate a bordo dell'Olandese sparire in una palla di
fuoco di antimateria all'interno del gigante gassoso. Un punto per me.

\hfill\break

Dopo aver visto l'esercitazione dell'equipaggio con armature potenziate
e combot, ero colpito, e iniziai a coltivare la speranza di poter
portare a termine l'assalto con successo. I nostri sei soldati
meccanizzati, con le loro tute spaziali kristang, erano rapidi e
potevano trasportare con facilità i pesanti fucili kristang. La gravità
dell'asteroide raggiungeva solo il 2\% della norma terrestre, quindi,
per l'addestramento, Skippy ridusse al 2\% la gravità nella stiva di
addestramento. Il 2\% significava che le persone in tuta spaziale e i
combot avevano bisogno di qualcosa di diverso dalla gravità per evitare
di staccarsi dal pavimento e galleggiare nella base di ricerca. I robot
si aggrappavano in automatico a pavimenti, pareti, soffitti, qualsiasi
cosa servisse loro per spostarsi dove l'operatore voleva che andassero:
avevano cuscinetti di presa su piedi, ginocchia, mani, gomiti, ovunque
potesse essere utile. Le tute kristang avevano pinze sotto le suole
degli stivali e sui palmi dei guanti, certo, ma anche su gomiti, spalle,
retro del casco, ginocchia e sul sedere. I nostri soldati in tuta
dovevano imparare ad accendere e spegnere e regolare la forza delle
pinze.

Nel complesso, a preoccuparci non era la capacità dei combot di eseguire
le manovre richieste. Erano le armature potenziate, che richiedevano
molta pratica per abituarsi. Il problema che i soldati in armatura
incontravano, nel labirinto di addestramento alla battaglia che Adams
aveva creato, non era girare gli angoli e attraversare le porte, era
ridurre la loro velocità e potenza. Si muovevano troppo in fretta e
rimbalzavano sui muri. Tre tute si ruppero durante l'addestramento, ne
avevamo sei rimaste a bordo della Fiore insieme a dei pezzi di ricambio,
quello che non avevamo erano molti umani alti più di un metro e ottanta.
Consigliai a Giraud di rallentare e andarci piano, non potevamo
permetterci che qualcuno si facesse male durante l'addestramento. Le
persone che indossavano le tute erano già abbastanza doloranti senza
sbattere contro le cose: anche sfruttando al massimo la capacità di
regolazione, le tute erano troppo grandi. Mettere dell'imbottitura negli
stivali per sollevare i piedi di chi li indossava fu di qualche aiuto.
Io stesso provai una tuta e, regolandola su un metro e ottanta, me la
sentii affondare nelle ginocchia, nell'inguine, sotto le braccia, e
dovetti allungare il collo per evitare di sbattere il mento sul fondo
del casco.

Le tute kristang erano impressionanti. I combot thuranin erano
grandiosi. C'erano due tipi di combot, noi usammo soltanto quello più
piccolo, perché Skippy aveva detto che il modello pesante rischiava di
essere troppo grande per i passaggi da attraversare nella base
sull'asteroide. Il combot piccolo aveva tre gambe per la stabilità, e
tre braccia. Di queste, due servivano per la stabilità o l'arrampicata,
il braccio al centro serviva per gli attacchi delle armi, e Skippy aveva
fabbricato occhiali che permettevano ai nostri operatori sia di vedere
attraverso i sensori ottici del combot, sia di controllare le armi.
Queste ultime s'indirizzavano ovunque l'operatore guardasse, con
reticoli di puntamento negli occhiali per mirare. Il piano originale di
Skippy prevedeva che l'operatore controllasse l'innesco dell'arma
sbattendo le palpebre; Giraud lo aveva cestinato facendo notare che la
maggior parte degli ammiccamenti degli occhi umani sono involontari. Su
suggerimento di Giraud, il grilletto era controllato dal dito indice
dell'operatore, a sinistra o a destra, a seconda che la persona fosse
mancina o destrorsa. La scelta dell'arma, il fucile o il
lanciarazzi/granate, era controllata dal pollice dell'operatore. A bordo
dell'Olandese, non potevamo usare vere munizioni, e questo era un
problema, perché gli operatori non potevano rendersi davvero conto
dell'efficacia della loro arma. Skippy avvisò che, anche all'interno
della base, dove non era alto il rischio di aprire un buco nella spessa
scorza dell'asteroide e causare una breccia nel contenimento
atmosferico, avremmo dovuto usare i razzi e le granate con parsimonia
mentre entravamo. Tornando indietro, per coprirci, avremmo potuto
sparare a qualsiasi cosa vedessimo. Il piano era di far saltare in aria
l'intero asteroide in ogni caso, per cancellare qualunque traccia del
fatto che gli umani fossero stati lì.

\hfill\break

L'assalto vero e proprio andò perfettamente secondo il piano di Skippy e
Giraud. All'inizio. E poi tutto andò a puttane.

Prima che facessimo il salto nel sistema stellare che conteneva la base
sull'asteroide, l'intero equipaggio, a parte Desai e Walorski, si riunì
nell'hangar d'atterraggio. Era il momento di un ultimo controllo
dell'equipaggiamento, per assicurarci che tutto quello di cui il gruppo
d'assalto aveva bisogno fosse caricato a bordo delle due navicelle e
ogni cosa funzionasse a dovere. Un combot di riserva fu caricato a bordo
di ciascuna delle navicelle, rendendone il già stretto interno ancora
più angusto. Se i thuranin avevano costruito soffitti alti nelle loro
navi madri per accogliere gli ospiti, non avevano fatto altrettanto con
le loro navicelle. Erano così limitate negli spazi disponibili che le
sei persone destinate a indossare l'armatura kristang dovettero farlo
prima di salire a bordo, anche se in tal modo avrebbero prolungato il
tempo di permanenza nelle scomode tute.

Dopo avere supervisionato l'ispezione, Giraud tenne un discorso breve ma
efficace, molto meglio di qualsiasi cosa avrei potuto dire io. Non usava
un linguaggio impetuoso, né cercava di ispirare le persone, non avevano
bisogno d'ispirazione. Si limitò a dichiarare che avremmo avuto una sola
possibilità in quell'assalto, una sola possibilità di liberare la Terra
dai kristang. Costasse quel che costasse, non avremmo lasciato
l'asteroide senza il modulo di controllo del wormhole. Tutti avevano
studiato la riproduzione che Skippy ne aveva fabbricato: il modulo era
una scatola rettangolare sottile, lunga circa un metro e ottanta, e
larga quindici centimetri; Skippy disse che si piegava ed espandeva per
formare una X di tre metri d'altezza. Chiunque avesse visto un modulo,
lo avrebbe preso, se aveva indosso una tuta kristang; l'avrebbe
segnalato a chi la indossava e aspettato, se stava manovrando un combot.
Non potevamo rischiare di danneggiare il modulo tentando di prenderlo
con un combot, la cui immensa forza, anche nel caso di operatori più
delicati, schiacciava di regola le cose. Senza cibernetica, l'operatore
non aveva feedback a sufficienza per evitare di stritolare gli oggetti
nella presa degli artigli. I combot avrebbero fatto strada attraverso la
base, seguiti dalla retroguardia dei nostri uomini con indosso le tute.

Giraud mi disse che avrebbe preferito fare altre due settimane di
pratica di combattimento. Non lo disse all'equipaggio. Skippy era
ansioso di colpire la base sull'asteroide il prima possibile: vista la
situazione militare fluida nel settore, era preoccupato del fatto che i
kristang avrebbero rafforzato la loro base di ricerca sull'asteroide, o,
peggio, avrebbero impacchettato qualsiasi cosa di valore per portarla
via, disperdendo le cose tra le stelle. Non potevamo correre questo
rischio, perciò decisi che saremmo saltati in battaglia il prima
possibile.

Prima che lanciassimo l'assalto, Chang e Giraud ebbero una conversazione
con il sottoscritto. Volevano essere rassicurati sul fatto che, nel caso
in cui quell'assalto sembrasse sul punto di fallire e ci fosse il
rischio che il coinvolgimento dell'umanità venisse scoperto, non avrei
esitato a far esplodere l'asteroide e saltare via. Se non fossimo
riusciti a recuperare la Fiore, avrei dovuto vaporizzare anche quella
nave, era cosparsa di Dna umano. Chang e Giraud mi fecero promettere --
perché sapevano che da soldato questo significava andare contro ogni mio
istinto -- che non avrei esitato ad abbandonare i miei predoni e
fuggire. Se perdi una battaglia puoi ancora vincere la guerra.
Nonostante Giraud avesse detto all'equipaggio che l'assalto era la
nostra unica possibilità, l'umanità ne aveva ancora qualcuna, finché
Skippy era al sicuro e dalla nostra parte. Se l'assalto non fosse
riuscito, l'Olandese avrebbe potuto arruolare nuovi membri per
l'equipaggio su Paradiso o sulla Terra e continuare a cercare il modulo
di controllo del wormhole; non mi veniva in mente come la cosa avrebbe
potuto funzionare, ma non era del tutto impossibile.

Lo promisi. Mi sentivo una merda, ma lo feci. Lo dovevo al nostro
equipaggio. Loro erano disposti a rischiare la vita in quella folle
missione, io dovevo capire che la missione valeva la loro vita, comunque
andasse a finire.

Walorski non aveva idea di quanto fossi restio a rimanere a bordo
dell'Olandese mentre il nostro equipaggio, il mio equipaggio, partiva
per l'assalto.

\hfill\break

Saltammo vicino all'asteroide, nel raggio di circa tre milioni di
chilometri. Il salto e le manovre successive erano stati pre-programmati
da Skippy, perché i tempi erano molto stretti; Desai non dovette fare
altro che premere un pulsante al mio comando. Una frazione di secondo
dopo il salto, la Fiore fu espulsa, eseguì un micro-salto e l'Olandese
fece un salto a corto raggio nella direzione opposta. La Fiore formò un
punto di salto per un altro micro-salto, poi il campo collassò come
pianificato e la Fiore cadde, morta nello spazio. L'effetto, secondo la
logica di Skippy, era che l'area circostante fosse sommersa dalle onde
sovrapposte del campo di salto, mascherando la presenza di una seconda
nave. La Fiore ruzzolò fuori controllo, i propulsori macinavano a caso,
la nave perdeva radiazione. Volevamo che attirasse l'attenzione, mentre
l'Olandese sfrecciava via, non vista. Skippy indusse la Fiore a
trasmettere i codici Iff di una fregata scomparsa nella zona due anni
prima, durante una schermaglia con un altro clan kristang: la nostra
speranza era che le lucertole sarebbero state abbastanza confuse e
curiose da non far saltare subito la Fiore in mille pezzi. Più i
kristang avrebbero indagato su di essa, meno avrebbero notato
l'Olandese.

Funzionò. Prima che i kristang dirigessero i loro sensori in modo da
allargare il raggio di scansione, Skippy era nella loro rete e istruiva
il sistema perché ignorasse l'enorme nave madre thuranin che si
avvicinava furtiva alla loro porta d'ingresso. Mentre attraversavamo le
loro linee di campo, facemmo scattare i reticoli di rilevamento stealth
e i sensori ci identificarono, ma i computer kristang semplicemente
ignorarono quegli input. Desai fermò l'Olandese dietro un asteroide a
tremiladuecento chilometri dal bersaglio e Skippy lanciò le navicelle.

La prima volta che Skippy ci aveva detto che avremmo fatto un assalto su
un asteroide, l'immagine che avevo in mente era un pezzo di roccia
irregolare e frastagliata, come gli asteroidi nei film spaziali. Questo
era tecnicamente un planetoide, aveva un diametro di
quattrocentottantadue chilometri ed era abbastanza grande, tanto che la
sua gravità l'aveva trasformato in una sfera. Eppure, era un ammasso di
roccia butterato, di un brutto colore grigio e marrone, gelido e
desolato. Da un lato c'era la base di ricerca, dall'altro una base
militare kristang molto più grande. Una grande base militare, con un
migliaio di soldati, navi d'assalto, cannoniere e abbastanza missili da
far saltare in aria l'Olandese Volante. Da qualche parte nella zona
c'erano uno squadrone di fregate, sei cacciatorpediniere e un
incrociatore. E stavamo per fare irruzione, prendere quello che volevamo
e andarcene via di colpo.

Skippy si dimostrò affidabile: la base militare era interessata alla
Fiore, e incaricò un paio di fregate e quattro navicelle di indagare
sulla nostra esca. La base militare inviò un messaggio alla base di
ricerca perché andasse in massima allerta; la base di ricerca non
ricevette mai il messaggio perché Skippy lo intercettò, e rispose per
loro, facendo credere alla base militare che la base di ricerca fosse
stata messa in sicurezza. Quando le navicelle si avvicinarono alle porte
dell'hangar d'attracco, li ponemmo di fronte a una quasi completa
sorpresa. Il sovrintendente dell'hangar pensava che le porte si stessero
aprendo per un paio di navicelle kristang inviate a rafforzare la
sicurezza, ed è quello che Skippy aveva detto ai computer di mostrare.
Le nostre navi atterrarono e schierarono i combot con i nostri sei
soldati in armatura potenziata. Tutto sarebbe andato alla perfezione,
non fosse che un tecnico della manutenzione kristang, in tuta spaziale,
stava lavorando al meccanismo appena oltre le porte. Vide con i suoi
occhi che le nostre navicelle erano thuranin, Skippy non poteva certo
hackerare il suo nervo ottico, e questo lavoratore devoto fece il suo
dovere e gridò un avvertimento alla radio. Un avvertimento sepolto da
Skippy. Poiché nessuno rispondeva, lo sgobbone kristang attraversò
l'hangar e abbassò una leva che fece scattare un allarme. L'unico modo
che aveva Skippy per sopprimere un allarme cablato era bruciare il
circuito elettrico, generando un sovraccarico di corrente dopo che
l'allarme però era squillato due volte. Skippy ordinò al computer di
annunciare che l'allarme era falso, che era stato innescato dal
sovraccarico.

Questo ci avrebbe fatto guadagnare altri due minuti, tempo a sufficienza
perché la squadra d'assalto potesse piazzare cariche esplosive sulla
porta interna e farla saltare in aria in modo controllato. L'avrebbe
anche fatto, ma il cittadino modello kristang capì che c'era qualcosa
che non andava quando l'allarme suonò soltanto due volte e si precipitò
verso un armadietto delle armi. Questo attirò l'attenzione del gruppo
d'assalto, che diede al nostro impiegato dell'anno kristang la sua
ricompensa, non sotto forma di denaro o posto auto privilegiato, ma di
fuoco concentrato proveniente da una mezza dozzina di combot. Ed è lì
che iniziarono i problemi, quando il kristang scomparve in una nuvola di
vapore insanguinata. I proiettili esplosivi che i nostri tesi pirati
avevano selezionato per i loro combot passarono il kristang da parte a
parte e colpirono il muro dell'hangar d'attracco. I kristang nella base
di ricerca non ebbero bisogno che un mezzo elettronico notificasse loro
che c'era qualcosa che non andava, poterono sentire la vibrazione
provocata dai proiettili esplosivi attraverso pavimenti e pareti.

A quel punto, Giraud ordinò ai combot di sgomberare il passaggio e
distrusse la porta interna con un razzo, con la testata impostata su
carica cava. La porta esplose verso l'interno, cosa per nulla ottimale,
perché parte di essa ciondolava ancora in pezzi frastagliati appesi al
telaio. Fu il momento in cui la nostra mancanza di esperienza si ritorse
contro di noi. Giraud avrebbe dovuto programmare la testata del razzo in
modo da ottenere un'esplosione a più ampio raggio, invece che alla
massima penetrazione: la porta era dura, ma non blindata. A quel punto,
far saltare il muro su entrambi i lati della porta per liberare un
percorso avrebbe fatto volare detriti dappertutto, perciò perdemmo tempo
manovrando un paio di combot perché strappassero i resti della porta con
le pinze. Uno di questi era un po' troppo entusiasta e usò eccessiva
forza, un grosso pezzo di porta e telaio attraversò in volo l'hangar
d'attracco nella bassa gravità e mancò di poco un altro combot. Il suo
operatore umano non ebbe il tempo di reagire: i sistemi automatici del
combot fecero in modo che la macchina lo schivasse, permettendo però al
pezzo di porta di colpire Chang sul lato sinistro. L'impatto lo fece
volare attraverso l'hangar, ruppe la tuta pressurizzata e gli incrinò
diverse costole. In seguito, scoprimmo che aveva due costole rotte. Una
volta ero caduto da una bici da cross e mi ero rotto una costola, il
dolore mi aveva fatto strisciare sulle ginocchia, cercando di riprendere
fiato, non fosse che respirare intensificava le fitte. Avevo dovuto
trascorrere un'ora seduto o sdraiato a terra prima di poter risalire
sulla bici e tornare, a ritmo molto lento, a casa, da dove i miei
genitori mi avevano portato all'ospedale. Il rivestimento della tuta
kristang di Chang conteneva un gel che s'induriva con l'esposizione al
vuoto, sigillando il buco per evitare che ne fuoriuscisse l'aria. In
qualche modo, anche se doveva provare un dolore incredibile, Chang si
rimise in piedi e di nuovo in azione, grazie a fegato e adrenalina. Il
ragazzo sputava sangue, ma questo non lo fermava. Skippy ci mostrò i
dati ricavati dai rilevatori medici integrati nella tuta di Chang:
sembrava grave, non capivo come facesse a non essere raggomitolato sul
ponte, e avrei voluto ordinargli di ritirarsi su una navicella e
lasciare che la squadra di assalto continuasse la missione. Non lo feci.
Avevo bisogno di credere che, se Chang non fosse stato in grado di
andare avanti, me lo avrebbe detto. E poi Giraud era lì con lui, se
avesse pensato che sarebbe stato inefficace in combattimento, me lo
avrebbe segnalato in privato. Avevamo bisogno di Chang, avevamo bisogno
di tutti. Finché non fosse stato d'ostacolo alla missione, non avrei
interferito da sedicimila chilometri di distanza. I motori della sua
armatura potenziata e la bassa gravità potevano essere stati d'aiuto,
assorbendo parte della tensione d'urto, bisognava vedere come avrebbe
gestito il combattimento.

Una volta sgombrata la porta interna, il gruppo d'assalto si precipitò
lungo il corridoio, davanti i combot, a seguire le sei persone in tuta,
infine due combot in retroguardia. Presto s'imbatterono in un'altra
porta, che impiegarono trenta secondi a far saltare in aria con degli
esplosivi. Dopo di che, si trovarono nella parte principale della base
di ricerca, con Skippy che poteva aprire e chiudere le porte per loro.

Riscontrammo subito un problema. A quel punto Skippy si era infiltrato
nel computer della base e aveva perquisito gli archivi che, secondo lui,
erano un miscuglio osceno di spazzatura ammucchiata a caso, non
indicizzata e non catalogata. A dispetto delle difficoltà, localizzò i
due oggetti che volevamo rubare, un modulo di controllo del wormhole e
una sorta di nodo di comunicazione associato al Collettivo. Erano
depositati in diversi compartimenti, per fortuna nessuno dei due era
considerato una priorità dai ricercatori kristang e si trovavano
entrambi in aree di bassa sicurezza, lontano dalle principali strutture
di ricerca. Dovetti decidere in fretta se ordinare al gruppo d'assalto
di andare a cercare gli oggetti uno dopo l'altro, oppure dividersi.
Dividere una forza andava contro i principi di guerra che l'esercito mi
aveva insegnato. Mentre ci stavo pensando, mi passavano per la testa le
parole del Manuale di campo dell'esercito americano. Dividere una forza,
in una situazione in cui ci trovavamo già in condizione d'inferiorità
per numero e armamenti, violava i principi militari della massa e
dell'economia della forza. Non aveva importanza, la possibilità di dover
dividere il gruppo d'assalto era stata discussa in anticipo e Chang,
Giraud, Simms e io c'eravamo trovati d'accordo. In questo caso, la
concentrazione della forza era secondaria al principio della sorpresa.
Il nostro gruppo d'assalto poteva avere successo, poteva sperare di
sopravvivere solo entrando e uscendo il più velocemente possibile, prima
che i kristang potessero riprendersi dalla sorpresa, capire cosa diavolo
stava succedendo e concentrare le proprie forze contro di noi. Più a
lungo il gruppo restava nella base, più a lungo era esposto al pericolo.
Diedi ordine a Chang di guidare un gruppo alla ricerca del nodo di
comunicazione e a Giraud di guidarne un altro alla ricerca del modulo di
controllo del wormhole. Quello che l'ordine non diceva era che sapevo
che Chang, ferito, era il più debole dei due, e io desideravo il modulo
di controllo del wormhole molto più di quanto non desiderassi una radio
sofisticata per consentire a Skippy di parlare con le antiche
intelligenze artificiali. Per un momento, pensai di tenere un paio di
combot di guardia nel punto in cui i gruppi si erano divisi, per
sorvegliare un percorso utile alla loro ritirata, finché Skippy non mi
mostrò una pianta della base: c'erano un sacco di incroci di corridoi
attraverso i quali i due gruppi sarebbero dovuti transitare,
sorvegliarne solo uno era uno spreco di potenza di fuoco.

Da quel momento, l'incursione procedette in fretta e in modo caotico,
per la maggior parte secondo i piani. Le lucertole erano molto
destabilizzate, avevano rafforzato i loro sistemi interni in previsione
di un attacco da parte dei clan rivali kristang o thuranin, ma non
avevano previsto che Skippy s'infiltrasse da cima a fondo nei loro
sistemi. I sistemi dei quali non aveva il controllo, li aveva
indeboliti, confusi, cortocircuitati o privati di energia. I kristang
avevano inviato frenetici appelli alla base militare dall'altro lato
dell'asteroide, Skippy aveva intercettato i messaggi e aveva inviato
risposte false, comunicando che la base militare era sotto attacco
pesante e avrebbe inviato rinforzi quando sarebbe stato possibile. Un
certo kristang dal pensiero super-veloce, senza dubbio secondo
classificato al concorso da dipendente dell'anno, sparò un razzo di
segnalazione che avrebbe dovuto comunicare alla base militare che
qualcosa andava molto male. Il razzo fece un grande fuoco d'artificio
quando passò all'orizzonte ed ebbi un momento di panico. Skippy mi calmò
assicurandomi che, anche se la rete di sensori kristang aveva percepito
la segnalazione luminosa, lui aveva fatto in modo che la ignorasse. La
base militare era sepolta sotto la roccia, a meno che non ci fossero dei
kristang in superficie per qualche motivo, nessuno avrebbe visto il
razzo e, dato che Skippy stava disturbando il traffico radio, se anche
ci fosse stato un kristang in superficie, avrebbe dovuto vedere il
razzo, sapere cosa significava, correre alla base militare e bussare
alla porta. Pensai che il rischio era basso.

Skippy prevenne il lancio di altri razzi, incendiando il contenitore in
cui erano custoditi, senza aprire la porta del silo. L'esplosione che ne
derivò potrebbe aver rovinato le prospettive del secondo classificato al
concorso per l'impiegato dell'anno di ricevere una lucida targa da
mettere sulla sua scrivania, dal momento che spazzò via un pezzo
piuttosto importante della base e con molta probabilità lo uccise. Per
quanto mi riguarda, non avevo nulla in contrario.

I combot stavano avendo successo, fu d'aiuto il fatto che la forza di
sicurezza kristang, secondo il protocollo stabilito, si era ritirata per
proteggere il centro di ricerca di alta sicurezza della base,
l'obiettivo più ovvio, e nessuno si erano accorto che miravamo a
un'altra parte della base. Skippy riferì che i kristang erano del tutto
confusi, presi dal panico, nel caos, incapaci di capire cosa stesse
succedendo e il motivo per cui i loro sistemi, rafforzati con cura, non
stessero funzionando e stessero lavorando addirittura contro di loro. Le
sovratensioni avevano bruciato i comandi di porte, ascensori, luci e
salvavita nelle parti della base di cui non ci importava; molti dei
kristang erano isolati, intrappolati, lenti a reagire e senza guida. La
poca resistenza che la squadra d'assalto incontrò durante il tragitto
era costituita da individui solitari o piccoli gruppi di due o tre, non
coordinati, inefficaci. Gli sforzi inutili compiuti dai kristang contro
di noi aiutarono di fatto i nostri operatori di combot ad acquisire
esperienza con il fuoco vivo, imparando in un lampo a usare raffiche
rapide, con colpi impostati sulla minima forza esplosiva e in modalità
frammentazione. Il rock`n'roll altamente esplosivo, usato dagli
operatori contro la prima coppia di kristang, fece a pezzi il nemico, ma
sprecò anche munizioni e fece saltare in aria il corridoio, creando
detriti che costrinsero la squadra d'assalto a manovre di aggiramento.
Giraud s'incazzò a ragione, e avvertì la gente di non dare di matto con
le armi thuranin avanzate. Dopo il primo incontro, iniziarono a prestare
attenzione.

La squadra di Chang fu la prima a raggiungere il suo obiettivo, ma trovò
soltanto una stanza disorganizzata piena di cianfrusaglie. Non
cianfrusaglie impilate su scaffali, etichettate nel modo opportuno,
nemmeno cianfrusaglie stipate in scatole. Cianfrusaglie ammucchiate a
caso, come se i kristang avessero aperto la porta, avessero gettato roba
dentro il locale e l'avessero richiusa.

«Oh, merda», gemette Skippy. «Non l'avevo previsto. Il database mostra
che il nodo di comunicazione è stato localizzato per l'ultima volta in
quella stanza, non ci sono sensori lì dentro, nemmeno una telecamera.»

«C'è qualche magia che puoi fare?», chiesi con ansia.

«Sto facendo del mio meglio», replicò Skippy sulla difensiva. «Per
questioni di segretezza, non ci sono molti sensori all'interno della
base di ricerca, a eccezione dei punti d'accesso e degli alloggi.
Inoltre, la base è stata progettata proprio per rendere difficile la
sorveglianza da remoto. Sono per lo più cieco lì dentro, non quasi cieco
come si aspetterebbero i kristang, ma abbastanza da mettermi a disagio.
C'è un buon numero di elementi della security kristang di cui ho perso
traccia, non so dove siano o cosa stiano facendo. Attraverso le
telecamere della squadra d'assalto, ho visto anche aree della base che
non sono negli schemi cui ho accesso, sto raccogliendo dati per
sviluppare un vero e proprio layout della struttura. Il mio consiglio
più sincero è che il gruppo d'assalto si muova il più in fretta
possibile.»

Wow, che grande idea, non ci avevo pensato. Questo è quello che mi passò
per la mente, non quello che dissi a Skippy. «Colonnello Chang», dissi
allo zPhone, «dovrete scavare in quel mucchio per trovare il nodo di
comunicazione. Sai com'è fatto.» Ogni membro della squadra d'assalto
aveva visto le riproduzioni dei due oggetti che dovevamo rubare
fabbricate da Skippy.

«Sì, ricevuto», disse Chang, la sua voce era tesa, aveva un rantolo in
gola.

Sugli schermi, potei vedere Chang, Darzi e Asok Putri entrare nella
stanza e iniziare la cernita nel mucchio. I movimenti del tenente
colonnello erano rigidi, a causa della ferita non scese sul pavimento,
ma iniziò dalla parte superiore del mucchio. Organizzava il lavoro,
sistemando in una pila gli oggetti che gli altri avevano già smistato,
in modo che non si confondessero. Erano efficienti, concentrati, più
rapidi possibile, eppure troppo, troppo lenti. Il gruppo di Chang era
molto in ritardo, la stanza era grande, con cinque mucchi di
cianfrusaglie alti fino al soffitto e, dopo cinque minuti, avevano
sistemato solo la metà di un mucchio. Anche la squadra di Giraud aveva
raggiunto il proprio obiettivo: era il suo giorno fortunato, perché la
stanza dove si trovava il modulo di controllo del wormhole era tanto in
ordine quanto il bersaglio di Chang era l'opposto. Giraud, il sergente
Thompson e il soldato Marsden corsero lungo i corridoi, mostrando il
contenuto degli scaffali a Skippy, che decifrò in fretta il sistema di
indicizzazione e diresse Thompson proprio dove si trovava il modulo di
controllo del wormhole, appoggiato su uno scaffale.

«Obiettivo raggiunto», riferì Giraud. «Torniamo indietro ora. Dobbiamo
tornare direttamente all'hangar d'attracco o aiutare la squadra del
colonnello Chang?»

La squadra di Chang stava impiegando troppo tempo. Merda, si aspettavano
che facessi questa chiamata, mandando la squadra di Giraud in aiuto di
quella di Chang, nella speranza di velocizzare le ricerche? Deviare la
squadra di Giraud rischiava di esporla più a lungo al pericolo e di
farci perdere il modulo di controllo del wormhole. Presi in
considerazione l'idea di abbandonare la ricerca del nodo di
comunicazione di Skippy, ora che avevamo il modulo di controllo del
wormhole? Non proprio. Avevo rimuginato sulla questione nell'ultima
settimana, mi aveva tenuto sveglio la notte. Ci interessava chiudere il
wormhole, era l'unica cosa che mi importava, era l'intero scopo della
missione, avrei sacrificato qualsiasi cosa per raggiungerlo. A noi umani
non importava davvero di Skippy. A me un po' importava, perché mi
piaceva, e sentivo che gli eravamo debitori, e a volte pensavo di aver
capito la sua terribile, antica, dolorosa solitudine. I miei sentimenti
personali dovevano essere messi da parte, però, in quanto comandante
avevo bisogno di concentrarmi sulla missione e sui miei uomini. Quello
che avevo deciso in anticipo era che, a prescindere dalla promessa che
avevo fatto a Skippy, se ci fosse stato un modo per chiudere il wormhole
senza di lui, non avrei messo a rischio la vita dei miei uomini e la
missione per prendere la sua radio magica. Ma non avevamo modo di
chiudere il wormhole da soli, il che rendeva facile la mia decisione,
anche se difficile da digerire.

«Tenente Giraud, porti la sua squadra in supporto del colonnello Chang,
dobbiamo muoverci il più in fretta possibile.»

Qualunque cosa pensasse Giraud della saggezza della mia decisione, non
discusse. «Ricevuto. Stiamo andando.»

Sullo schermo, vedevo che Chang, Putri e Darzi stavano smistando i
mucchi di cianfrusaglie il più in fretta possibile; Chang aveva ordinato
a un combot di entrare nella stanza senza toccare nulla, per utilizzarne
le telecamere, nella speranza che Skippy riuscisse a trovare la cosa. La
ricerca si stava protraendo troppo a lungo. Skippy mi disse che aveva
perso le tracce dei kristang, che anche lui era per lo più cieco,
mettendomi una grande ansia. «Skippy, questa cosa del nodo di
comunicazione serve per le comunicazioni, giusto?»

«Sì, ma tu pensa. Quindi?»

«Quindi puoi contattarlo, usarlo per inviare un segnale che ti può
servire a trovare quella cazzo di cosa?»

«Merda. Sì.»

«Ma tu pensa.» Non riuscii a resistere, in parte per via dei nervi tesi.

«Cazzo, a volte non so cosa c'è che non va in me. Certo, sì. Fai
allontanare Chang, Putri e Darzi l'uno dall'altro, così posso usare le
loro radio per triangolare.»

Così feci. Meno di un minuto dopo, Skippy lo localizzò, sul retro del
mucchio di cianfrusaglie. Chang disse a Darzi e Putri di scavalcare il
mucchio nel mezzo e buttare la merda inutile fuori dai piedi, finché non
ebbero scoperchiato il primo terzo del mucchio, poi procedettero con più
attenzione. Fu Darzi a trovare la stupida cosa, era proprio dove aveva
detto Skippy. Sullo schermo, vidi che la squadra di Giraud era vicina
alla posizione di Chang, abbastanza vicina da farmi ritenere che fosse
meglio per loro collegarsi e sostenersi a vicenda sulla via del ritorno,
anche se questo significava puntare tutto su una sola carta.

Fu allora che le cose andarono in fretta a puttane.

La prima avvisaglia di guai non fu un avvertimento da parte di Skippy,
non furono i combot che incontravano i combattenti kristang davanti a
loro. Fu una pioggia di proiettili che colpì Giraud e Putri alle spalle,
abbattendoli entrambi. Putri subì numerosi colpi centratissimi e morì
all'istante, per fortuna i proiettili con punta esplosiva mancarono il
modulo di controllo del wormhole, che si trovava in un'imbracatura sulla
sua schiena. Giraud cadde sotto un colpo di striscio al lato del casco,
un duro colpo alla parte inferiore della schiena e un proiettile che gli
colpì il braccio sinistro e gli tagliò l'avambraccio a metà tra polso e
gomito. Darzi reagì proprio come avrebbe dovuto, com'era stato
addestrato a fare, si lasciò cadere sul campo, spianando la strada per
consentire ai combot di entrare in azione. Strisciò in avanti tenendosi
basso e strappò le cinghie dell'imbragatura di Putri, vedendo che per
lui non c'era più niente da fare, e si concentrò nel modo opportuno sul
mantenimento dell'obiettivo della missione.

I nostri combot ruotarono per affrontare i sei soldati nemici che
indossavano pesanti armature, fu uno scontro a fuoco violento e Darzi
non poté fare nulla per aiutare, dovette tenere la testa bassa e
proteggere il modulo vitale dietro la sagoma immobile di Giraud. Sullo
schermo potevo vedere che Giraud era vivo, anche se incosciente, e non
sanguinava in modo serio.

«Skippy», gridai frenetico, «c'è un'altra strada per tornare all'hangar
d'attracco?»

Quattro dei kristang erano ancora vivi e si stavano riparando in accessi
e corridoi laterali, tenendo tutti i combot impegnati. Darzi non poteva
muoversi senza esporre il prezioso modulo. La squadra di Chang stava
correndo per dare manforte, ma non sarebbe arrivata prima di due minuti,
forse qualcosa meno. Comunque troppo.

«Sì, esistono molti percorsi alternativi.»

«Ottimo. Operatori combot, non preoccupatevi di danneggiare il
corridoio, torneremo per un'altra strada! Prendete quei kristang,
massima forza!»

Era l'unica cosa che gli operatori dei combot avevano bisogno di
sentire, tre di loro impostarono i razzi su massimo rendimento e carica
cava, e fecero scoppiare l'inferno nel corridoio oltre Darzi e Giraud. I
razzi penetrarono a fondo in due o tre compartimenti, distrussero i
quattro kristang, fecero esplodere muri, pavimenti e crollare il
soffitto.

«Darzi!», gridai quando la mia voce sarebbe dovuta restare calma.
«Alzati, il tenente Giraud è vivo. Riesci a trasportarlo?»

«Sì, signore.» Darzi poteva portare con facilità Giraud nella bassa
gravità, specie nella sua armatura potenziata. Si fece imbracare il
modulo sulla schiena, legò le cinghie in fretta e sollevò il corpo
inanimato di Giraud. I combot fecero loro strada; avvisai quando la
squadra di Chang fu vicina, in modo che i due gruppi non si sparassero
l'un l'altro. Si riunirono e Skippy li guidò su un percorso alternativo,
una via traversa che riportava all'hangar d'attracco; questa volta Chang
aveva un paio di combot a coprirgli le spalle, in modo da evitare brutte
sorprese. Io ero preoccupato, perché il percorso tortuoso avrebbe
richiesto più tempo; Skippy mi assicurò che la via che il gruppo
d'assalto stava prendendo avrebbe confuso il nemico, che non avrebbe
saputo dove andare. La mia intelligenza artificiale però era ancora per
lo più cieca lì dentro. Questo era il punto più pericoloso della
missione, prima il nemico doveva indovinare gli obiettivi della squadra
d'assalto, ora il nemico sapeva dove la squadra d'assalto si stava
dirigendo: tornava all'hangar d'attracco. Gli uomini si trovavano a due
terzi della via del ritorno, secondo i calcoli di Skippy, quando si
presentarono dei problemi. Questa volta il nemico non apparve in
armatura, ma con i propri combot, che fecero saltare in aria un muro e
balzarono nel corridoio. I combot kristang erano pesanti, ingombranti e
lenti rispetto ai nostri modelli thuranin; erano dei facili bersagli.
Erano controllati da operatori molto più esperti, tuttavia, e avevano
tutti i vantaggi possibili. Tranne uno. Mentre lo scontro a fuoco
infuriava e la squadra d'assalto era per il momento bloccata, Skippy
gridò eccitato che i movimenti goffi dei nostri combot in realtà ci
stavano favorendo: i kristang si erano allenati per combattere i
movimenti super-veloci dei combot controllati dai thuranin e le lente
azioni impacciate dei nostri facevano loro sbagliare mira.

Tre combot li perdemmo nello scontro a fuoco iniziale, un altro era
danneggiato e incapace di usare le sue armi primarie. Chang vide che
Darzi era intralciato dal trasporto di Giraud e gli ordinò di metterlo
tra le braccia del combot danneggiato. Il soldato Marsden si fece avanti
per aiutare Darzi a legare Giraud, avevano quasi finito quando un
proiettile prese Marsden alla testa e gli frantumò la visiera, poi un
secondo quasi separò il casco dal collo. Cadde con un getto di sangue
che si congelò all'istante e sullo schermo vidi i suoi segnali vitali
spegnersi in una linea retta. Avvisai Chang di non perdere tempo ad
assistere Marsden, non si preoccupò di dirmi che aveva ricevuto, lo
sentii e lo vidi sollecitare la gente a proseguire. Ci potevano essere
altre strade per l'hangar d'attracco, e avevamo perso uomini avanzando
nello scontro a fuoco: Chang sapeva che era probabile che avremmo subito
ulteriori perdite se fossero rimasti bloccati sull'asteroide un secondo
più del necessario. Velocità e manovre erano fondamentali. Chang lo
vide, anche attraverso la dolorosa nebbia di costole rotte. Era un Iron
Man, con o senza armatura potenziata.

Dietro un muro di potenza di fuoco, la squadra d'assalto avanzava, con i
combot a fare strada. Una volta superato quel grande gruppo di combot
kristang, abbatterono la resistenza sparsa di altri combot e lucertole
in armatura. Iniziavo a preoccuparmi, eravamo partiti con sei persone in
tuta spaziale e nove combot; al momento di combattere per aprire un
varco verso l'hangar d'attracco, avevamo due morti, un ferito privo di
sensi e solo tre combot ancora funzionanti. Questi ultimi rimasero a
guardia della porta interna, mentre Chang prese il sergente Thompson e
Darzi a bordo di una navicella e portò Giraud a bordo dell'altra. Skippy
aveva già in parte retratto la porta esterna e, non appena Thompson e
Darzi ebbero allacciato le cinture, fece partire a manetta e decollare
quella navicella per uscire e andarsene via. Vidi sullo schermo che
spingeva a 6 G e gli ordinai di rallentare un po', una volta che la
navetta si trovò a sedici chilometri di distanza.

La navicella di Chang aveva accumulato un leggero ritardo mentre
mettevano le cinture di sicurezza a Giraud, e quel ritardo ci aveva
quasi ucciso. Thompson aveva il nodo e Darzi aveva il modulo:
tecnicamente avevamo già portato a termine la missione, non fosse per le
persone nella seconda navicella. E fu allora che un proiettile vagante
attraversò la porta interna, mancò i nostri combot, attraversò l'hangar
d'attracco e colpì lo scafo blindato della nave thuranin. La punta
esplosiva eruttò verso l'interno, creando un getto di plasma
surriscaldato, che attraversò bruciando la corazza fino all'interno
della navicella. Una volta che il plasma colpì l'aria, esplose in tutte
le direzioni. Il sergente scelto dell'Aeronautica statunitense Joy
Chung, che manovrava uno dei tre restanti combot, fu ucciso all'istante
e altri tre uomini riportarono gravi ferite. Chang e Giraud erano
protetti dalle loro armature, tutti gli altri furono colpiti duramente.
Con persone che urlavano di dolore, sangue che volava nella navicella e
aria che fuoriusciva dal buco nello scafo, Chang non esitò. Ordinò a
Skippy di farli uscire da quel cazzo di posto. La navicella sfrecciò
fuori dall'hangar d'attracco e si allontanò dalle porte, mentre Skippy
prendeva il controllo dei sistemi di riparazione automatizzati per
tappare il buco largo quanto un dito e far pompare di nuovo l'aria.

«Libere. Entrambe le navicelle libere», riferì Skippy in tono obiettivo.
«Nessun inseguimento. Pacchetti al sicuro. Sto guidando le navicelle
verso di noi.»

«Grande», dissi, emotivamente esausto. «Pilota, sarò nel nostro hangar
d'attracco.»

«No, signore», disse Desai piano. «Il colonnello Chang e il maggiore
Simms possono gestire le persone che hanno bisogno di aiuto. Di lei
abbiamo bisogno qui.»

Aveva ragione e lo sapevo. Avrei potuto farmi prendere dall'emotività
più tardi. Più tardi, dopo che saremmo saltati via in sicurezza.
«Capitano Desai, hai ragione», abbassai lo sguardo sulle mie mani
tremanti per il disgusto. Non ero stato in pericolo neanche per un
istante. Bisognava aspettare otto lunghi minuti perché le navicelle
tornassero.

\hfill\break

«Navicelle al sicuro. Pilota, iniziare salto, opzione Alfa», annunciò
Skippy.

«Ferma», ordinai, e Desai si voltò verso di me confusa. «Salta in
opzione Charlie, ora. Vai.»

Desai premette il pulsante giusto e la nave eseguì un micro-salto,
quindi ora ci trovavamo in uno spazio libero, non schermato da un
asteroide. Skippy aveva programmato cinque opzioni di salto nel sistema
di navigazione e io le avevo memorizzate, erano anche elencate sullo
schermo della poltrona di comando e sul mio iPad. L'opzione A prevedeva
un salto in prossimità della Fiore, così da poterla recuperare; il piano
prevedeva di saltare vicino al velivolo, inviare un segnale perché
quest'ultimo eseguisse un salto a corto raggio, far coincidere il salto
con l'Olandese, prendere la Fiore su un hangar d'attracco e poi saltare
nel bordo esterno del sistema stellare. L'opzione C prevedeva un
micro-salto lontano dall'asteroide dietro il quale ci nascondevamo, nel
caso in cui la nostra posizione fosse stata scoperta, ma dovessimo
rimanere nell'area per recuperare le navicelle.

«Posso chiedere perché non andiamo a prendere la Fiore?», chiese Skippy.

«Lo faremo. Segnala alla Fiore di saltare nel punto d'incontro.» Pur
strisciando alla velocità dalla luce, il segnale avrebbe raggiunto la
nostra fregata kristang rubata in pochi minuti e quella nave,
all'apparenza abbandonata, si sarebbe animata all'improvviso e avrebbe
acceso i motori per un salto a corto raggio.

«Fatto», riferì Skippy. «Colonnello Joe, perché siamo ancora qui?»

«Hai il controllo della testata nucleare vicino al loro hangar
d'attracco principale?»

«Sì, certo. Te l'ho detto.» I kristang avevano attrezzato la sezione di
ricerca della base sull'asteroide in modo che potesse autodistruggersi
con una testata nucleare; se qualcuno avesse invaso la base e i kristang
ne avessero perso il controllo, avrebbero potuto farla saltare in aria
con un'unica esplosione. La presenza della testata serviva a far sì che
i nemici decidessero che non valeva la pena assalire la base. Di certo,
avrebbe scoraggiato anche noi, se Skippy non ne avesse preso il
controllo poco prima che lanciassimo le navicelle.

«Quanto è grande?»

«Presumo tu stia chiedendo della sua resa esplosiva: quell'arma è di
circa otto megatoni. Una tipica testata W76, nei missili americani
Trident lanciati da sottomarino, è di cento chilotoni, per darti un
metro di paragone.»

«Innescala.»

«Cosa?» Skippy, Desai e Walorski ebbero la medesima reazione.

«Skippy, non discutere. So che hai cancellato la memoria del computer,
ma abbiamo lasciato del Dna lì dentro, sangue umano. Potrebbe rivelare
ai kristang che siamo coinvolti nell'assalto. Non posso correre questo
rischio. Innescala. Ora.»

A suo credito, va detto che Skippy sapeva quando non c'era da discutere:
«Arma abilitata sul grande pulsante rosso».

Premetti il pulsante. L'asteroide avvampò in una luce intensa. Desai ci
fece saltare via prima che i detriti ci colpissero.